% recap.tex
% Section 2: Background

\section{Background}
\label{sec:recap}

Equilibrium and incentive compatibility are both defined by deviations.
Computing an equilibrium, or verifying incentive compatibility for a candidate
mechanism, can be costly. Learning can share computation across related
instances, but it does not itself establish the relevant guarantee. A learned
output therefore still needs a check or bound specific to the claimed property.

\subsection{Equilibrium Conditions and Computational Demands}
\label{subsec:recap-gt}

For both concepts, the relevant deviation class determines what must be checked.
We first state the corresponding deviation conditions, then separate
constructing a candidate from verifying it. The final subsection reviews
standard restrictions and relaxations that recover tractability.

\subsubsection{Equilibrium and its deviation class}

Von Neumann's minimax theorem establishes that every finite two-player zero-sum
game has a value and a pair of optimal mixed
strategies~\cite{Neumann1928,vNM1944}. In that setting a player's equilibrium
strategy is also its worst-case-optimal strategy, so equilibrium and robustness
coincide. This coincidence is specific to strictly competitive play and does not
extend to general games.

Nash's equilibrium is a profile in which no player gains by deviating
alone~\cite{Nash1951}. The condition can be quantified through each player's
best unilateral gain. For a profile $s$ and a class $\mathcal{D}_i$ of
deviations available to player $i$, write the aggregate gap
\[
  \mathrm{Gap}_{\mathcal{D}}(s)
  \;=\;
  \sum_i \Bigl[\, \sup_{s_i' \in \mathcal{D}_i} u_i(s_i', s_{-i}) - u_i(s) \,\Bigr].
\]
Assume that $\mathcal{D}_i$ includes the player's current strategy. When each
$\mathcal{D}_i$ is the player's whole strategy set, this sum is commonly called
NashConv; closely related normalizations are called exploitability. It is zero
exactly at a Nash equilibrium. In other applications $\mathcal{D}_i$ may be a
restricted class, in which case the gap describes only deviations in that class.

The deviation inequalities can be evaluated one player at a time while holding
the opponents' strategies fixed, but every player's inequality must hold for
the same joint profile. The condition is static and does not describe whether an
adaptive process reaches such a profile. That requires a convergence analysis of
the process itself.

Extensive-form games represent the order of play and the information available
to each player, and Kuhn's
theorem established that under perfect recall mixed and behavioral strategies
induce the same outcome distributions~\cite{Kuhn1953}. Selten's refinements
imposed rationality at contingencies the equilibrium path never
reaches~\cite{Selten1975}. Harsanyi handled private information by making a
player's payoff-relevant characteristics a \emph{type}, converting incomplete
information into a Bayesian game~\cite{Harsanyi1968}. In each case the
comparison with counterfactual alternatives remains the content of the concept,
with the refinements changing the set of alternatives being compared.

In a correlated equilibrium, a mediator draws a joint action profile and
privately recommends one action to each player. No player gains by deviating as
a function of the recommendation received~\cite{Aumann1974}. Quantal response
equilibrium replaces exact best response with a smooth choice rule in which
better actions are more likely, but mistakes persist~\cite{McKelveyPalfrey1995}.

\subsubsection{Design and deviation-based incentive constraints}

Mechanism design chooses rules under which participants' incentives produce a
desired outcome.
Impossibility results identify restrictions on what particular classes of
rules can achieve~\cite{Arrow1951,Gibbard1973,Satterthwaite1975}. Under the
conditions of the revelation principle, the design problem can be posed over
direct mechanisms in which agents report types and truthful reporting is an
equilibrium.

Incentive compatibility is likewise defined by deviations: for every relevant
type, truthful reporting must do at least as well as \emph{every} admissible
alternative report. Here the deviation class is the misreport space, with its
scope determined by the equilibrium concept and the mechanism's domain. The
Vickrey--Clarke--Groves family achieves efficient allocation in dominant
strategies~\cite{Vickrey1961,Clarke1971,Groves1973}, and Myerson characterized
the revenue-optimal single-item Bayesian
auction~\cite{Myerson1981}.

\subsubsection{Construction and checking}

The PPAD-completeness results noted at the start of the introduction concern
total search problems: a solution exists, although finding one may be
intractable~\cite{Papadimitriou1994}. The asymmetry between construction and
checking also depends on the representation and the property. Given an explicit finite game and a rational profile,
evaluating a specified deviation is direct. Finding one profitable deviation
refutes the equilibrium claim. Establishing that no profitable deviation exists
over a large or implicitly represented class can require an optimization or
structural result. Mechanism design faces the
same computational pressure because optimization over contingent rules is
itself a hard search problem~\cite{ConitzerSandholm2002}.

\subsubsection{Standard routes to tractability}

Tractability can be recovered by weakening the solution concept or restricting
the game class, often while accepting a bounded residual gain. For example, correlated
and coarse correlated equilibria in explicitly represented finite games admit
linear descriptions and are computable in polynomial time. They can also arise
from decentralized learning: when every player in a repeated game
drives external regret to zero, the empirical distribution of play approaches
the coarse correlated set, and stronger regret notions reach the correlated
set~\cite{Hannan1957,HartMasColell2000,BlumMansour2007}. The deviation class is
explicit in this result: the equilibrium notion obtained depends on which
deviations the regret criterion compares against~\cite{GordonGreenwaldMarks2008}.
Each modification changes the claim for which a deductive guarantee is
available. Learning can be combined with any of them by fitting a conditional
rule across related games and reusing it on later instances.

\subsection{What Training Does}
\label{subsec:recap-dl}

Section~\ref{subsec:recap-dl} first explains how training amortizes computation
across instances, then turns to training cost and the guarantees that must still
be established after training.

\subsubsection{Amortization: sharing computation across instances}

A per-instance solver is invoked on a specified game. An amortized system fits
a conditional rule across a family of related cases and later evaluates that
rule on a new case. Total computation includes training and later evaluations.
The potential advantage comes from sharing the trained rule across cases. This
is the structure of amortized optimization and
inference~\cite{Amos2022Amortized,KingmaWelling2014}. Different architectures
encode different forms of sharing across states or instances.

Reinforcement learning adds the dependence that matters for strategic settings:
an action changes the data the learner will next
observe~\cite{SuttonBarto2018,Williams1992,SuttonPolicyGradient2000,Mnih2015}.
When the rules are known, self-play can generate new training games as the
policy changes, an approach used to train policy and value networks in board-game
systems~\cite{Silver2018}. When behavior is observed
but the objective is hidden, inverse reinforcement learning estimates which
objectives are consistent with that behavior~\cite{NgRussell2000,HoErmon2016}.

Preference-based fine-tuning turns
comparisons between outputs into a training signal, making the learning target
itself something another party
supplies~\cite{Christiano2017,Ouyang2022}. Large language models generate
general-purpose text~\cite{BrownGPT3_2020}. Code models and tool-using agents
also produce executable programs or action traces
~\cite{ChenCodex2021,YaoReAct2023,SchickToolformer2023}.

\subsubsection{What must still be accounted for}

Amortization changes where the computation occurs. Data generation and parameter
optimization can dominate the time saved at inference, so a cost comparison
should include training as well as later evaluations. The strategic property of
the output then depends on the analytic or post-training check used by the
method.

Validation also depends on access. A high-dimensional rule may be available only
through queries, and finite tests cover sampled instances and representable
alternatives. Extending such results to a larger deviation class requires a bound
or structural argument. Executable code and rules constrained to a known family
may instead permit symbolic checking or a direct theorem.

Construction and validation impose separate costs. Learning can amortize
construction across related instances, but it does not by itself establish the
target property over the full deviation class. Sections~\ref{sec:ai4gt}
and~\ref{sec:gt4ai} identify the additional argument required in each
application domain.
